\documentclass{article}
\usepackage[utf8]{inputenc}
\usepackage{geometry}
\geometry{
	a4paper,
	total={170mm,257mm},
	left=20mm,
	top=20mm,
}
\usepackage{graphicx}
\usepackage{longdivision}
\usepackage[dvipsnames]{xcolor}
\graphicspath{{images/}}
\usepackage{svg}
\usepackage{titling}
\usepackage[american,RPvoltages]{circuiTikz}
\usepackage{tikz, tikz-cd}
\usepackage{pgfplots}
\usepackage{siunitx}
\usepackage{float}
\usepackage{enumitem}
\usepackage{amsmath,amsfonts,stmaryrd,amssymb} % Math packages
\usepackage{tcolorbox}
\usepackage{pdfpages}
\usepackage{bodegraph}
\usepackage{gensymb}
\usepackage{amsmath}
\usepackage{scalerel}
\setcounter{MaxMatrixCols}{20}

\newcommand{\longdiv}{\smash{\mkern-0.43mu\vstretch{1.31}{\hstretch{.7}{)}}\mkern-5.2mu\vstretch{1.31}{\hstretch{.7}{)}}}}


\title{Average Value}
%\author{Rodrigo Castillejos}
\author{Belmont}
\date{Mayo 2025}


\usepackage{fancyhdr}
\fancypagestyle{plain}{%  the preset of fancyhdr 
	\fancyhf{} % clear all header and footer fields
	\fancyfoot[L]{\thedate}
	\fancyhead[L]{Average Value}
	\fancyhead[R]{\theauthor}
}
\makeatletter
\def\@maketitle{%
	\newpage
	\null
	\vskip 1em%
	\begin{center}%
		\let \footnote \thanks
		{\LARGE \@title \par}%
		\vskip 1em%
		%{\large \@date}%
	\end{center}%
	\par
	\vskip 1em}
\makeatother

\renewcommand{\figurename}{Figura}
%\usepackage[color=ForestGreen!15,angle=45]{draftwatermark}
%\SetWatermarkText{CONFIDENTIAL\\COPY FOR\\ POCA POCA}
%\SetWatermarkScale{2.5}

\usepackage{lipsum}  
%\usepackage{cmbright}
%\usepackage{mathpazo}
%\usepackage{ccfonts,eulervm}
 \usepackage[T1]{fontenc}
 \usepackage{ccfonts}
%\usepackage[math]{anttor}
%\usepackage{antpolt}
 %\usepackage{fouriernc}
%\usepackage[utopia]{mathdesign}
%---------------------------------------------------
\tikzset{flipflop myRS/.style={flipflop,
		 flipflop def={t1=R, t3=S, t6=Q, t4={\ctikztextnot{Q}}, tu={\ctikztextnot{RST}}, nu=1, n4=1 }}
	}
	
\tikzset{flipflop myFFRS/.style={flipflop,
		flipflop def={t1=R, t3=S, t6=Q, t4={\ctikztextnot{Q}}, n4=1 }}
}
%---------------------------------------------------
\setlength\parindent{0pt}

\usepackage{mathtools}
\DeclarePairedDelimiter\abs{\lvert}{\rvert}%
\DeclarePairedDelimiter\norm{\lVert}{\rVert}%

% Swap the definition of \abs* and \norm*, so that \abs
% and \norm resizes the size of the brackets, and the 
% starred version does not.
\makeatletter
\let\oldabs\abs
\def\abs{\@ifstar{\oldabs}{\oldabs*}}
%
\let\oldnorm\norm
\def\norm{\@ifstar{\oldnorm}{\oldnorm*}}
\makeatother


\begin{document}
	
\maketitle
	
\noindent\begin{tabular}{@{}ll}
Author & \theauthor\\
\end{tabular}
	
\begin{figure}[H]
	\centering	
	\begin{tikzpicture}[scale=0.9]
		\pgfplotsset{width=9cm,compat=1.18}
		\begin{axis}[xshift=11cm,
			yshift=-7cm,
			clip=false,
			axis y line=left,
			axis x line=middle, 
			xmin=0,xmax=3.5,
			ymin=0,ymax=1.5,
			xlabel=$t$,
			ylabel=$i(t)$,
			axis y line= center,
			xtick={0.5,3},
			xticklabels={$t_1=DT$,$T$},
			ytick={0.1,1},
			yticklabels={$B$,$A$},
			]
			\addplot+[thick,mark=none,const plot]
			coordinates {(0,0) (0,1) (0.5,0.1) (3,0.1) (3,1) (3.1,1)};
			%\draw[<-](0.5, 1.05)--(0.7, 1.2); %vo
			%\draw[<-](1.2, 0.1)--(1.4, 0.3); %vc
			%\node[anchor=west] at (axis cs:0.6,1.3){$v_O$};
		\end{axis}
		;\end{tikzpicture}
	\caption{}	
	\label{fig:figura2}
\end{figure}

\begin{equation}
	i(t) = \begin{cases}
		A, \quad 0 \leq t \leq t_1\\
		B, \quad t_1 \leq t \leq T
	\end{cases}\,.
\end{equation}

The average value is given by:

\begin{align}
	I_{avg} = \frac{1}{T}\int_{0}^{T} i(t) dt
\end{align}

substituting (1) into equation (2) we get:

\begin{align}
		I_{avg} &= \frac{1}{T} \left[ \int_{0}^{t_1} A \, dt + \int_{t_1}^{T} B \, dt \right] = \frac{1}{T} \left[  At \bigg|_{0}^{t_1} + Bt \bigg|_{t_1}^{T} \right] \nonumber \\
		&= \frac{1}{T} \left[ At_1 + B(T - t_1) \right]
\end{align}

If we assume that $t_1 = DT$ where $D = \frac{t_1}{T} $ then:

\begin{align}
	I_{avg} = \frac{1}{T} \left[ ADT + B(T - DT) \right] 
\end{align}

\begin{tcolorbox}[colback=red!5!white,colframe=red!75!black,arc=0mm,boxrule=0mm,width=(\linewidth+0.5cm)]
	\begin{equation}
		I_{avg} = AD + B(1 - D)
	\end{equation}
\end{tcolorbox}

If $A= \qty{20}{\milli \ampere}$, $B = \qty{60}{\mu \ampere}$, $t_1 = \qty{20}{\mu \second}$, $T = \qty{10}{\milli \second}$,  then

\begin{align}
	D = \frac{t_1}{T} = \frac{\qty{20}{\mu \second}}{\qty{10}{\milli \second}} = 2 \times 10^{-3} \\
	I_{avg} = (\qty{20}{\milli \ampere})(2 \times 10^{-3}) + (\qty{60}{\mu \ampere})(1 - 2 \times 10^{-3} ) = \qty{99.88}{\mu \ampere}
\end{align}

\begin{align}
	T &\geq \frac{t_1(A + 0.01B)}{0.01B} = t_1 \left( 100 \frac{A}{B} + 1 \right)\\
	T &\geq \qty{666.69}{\milli \second}
\end{align}

\begin{align}
	I_{avg} &\approx B(1 - D) \\
	& \approx B\left( 1 - \frac{t_1}{T} \right)
\end{align}

\begin{align}
	Vo = \frac{ \qty{10}{\milli \volt} }{  ^{\circ}C } \times T[^{\circ}C] + \qty{500}{\milli \volt}
\end{align}

\begin{align}
	T[^{\circ}C] = \frac{V_o - \qty{500}{\milli \volt}}{ \frac{ \qty{10}{\milli \volt} }{ ^{\circ}C} }
\end{align}

\begin{align}
		T[^{\circ}C] = \qty{21.92}{\degreeCelsius}
\end{align}


\[
\arraycolsep=1pt
\renewcommand\arraystretch{1.2}
\begin{array}{*1r @{\hskip\arraycolsep}c@{\hskip\arraycolsep} *{11}r}
	&          & 1 & + & x & + & 2x^2 & + & 3x^3 & + & 5x^4 & + & \dots  \\
	\cline{2-13}
	1-x-x^2 & \longdiv & 1 &   &   &   &      &   &      &   &      &   &        \\
	&          & 1 & - & x & - &  x^2 &   &      &   &      &   &        \\
	\cline{3-7}
	&          &   &   & x & + &  x^2 &   &      &   &      &   &        \\
	&          &   &   & x & - &  x^2 & - &  x^3 &   &      &   &        \\
	\cline{5-9}
	&          &   &   &   &   & 2x^2 & - &  x^3 &   &      &   &        \\
	&          &   &   &   &   & 2x^2 & - & 2x^3 & - & 2x^4 &   &        \\
	\cline{7-11}
	&          &   &   &   &   &      &   & 3x^3 & + & 2x^4 &   &        \\
	&          &   &   &   &   &      &   & 3x^3 & - & 3x^4 & - & 3x^5   \\
	\cline{9-13}
	&          &   &   &   &   &      &   &      &   & 5x^4 & + & 3x^5   \\
	&          &   &   &   &   &      &   &      &   &      &   & \vdots \\
\end{array}
\]

\newpage

\begin{align}
	D_I = T - 0.55(1 - 0.01RH)(T-14.50) \nonumber  \\
	D_I - T = -0.55(1 - 0.01RH)(T-14.50) \nonumber \\
	T - D_I = 0.55(1 - 0.01RH)(T-14.50) \nonumber \\
	\frac{T - D_I}{T - 14.50} = 0.55(1 - 0.01RH) \nonumber \\
	1 + \frac{14.50 - D_I}{T - 14.50}  = 0.55(1 - 0.01RH) \nonumber \\
	\frac{14.50 - D_I}{T - 14.50} = 0.55(1 - 0.01RH) - 1 \nonumber \\
	\left[ 0.55(1 - 0.01RH) - 1 \right] \left[ T - 14.50 \right] = 14.50 - D_I \nonumber \\
	T - 14.50 = \frac{14.50 - D_I}{0.55(1 - 0.01RH) - 1} \nonumber
\end{align}

\begin{tcolorbox}[colback=red!5!white,colframe=red!75!black,arc=0mm,boxrule=0mm,width=(\linewidth+0.5cm)]
	\begin{equation}
		T = 14.50 + \frac{14.50 - D_I}{0.55(1 - 0.01RH) - 1} \nonumber
	\end{equation}
\end{tcolorbox}

Where:

\[
\arraycolsep=1pt
\renewcommand\arraystretch{1.2}
\begin{array}{*1r @{\hskip\arraycolsep}c@{\hskip\arraycolsep} *{11}r}
	&          &  &  & 1 &  &  &  &  &  &  &  & \\
	\cline{2-8}
	T-14.50 & \longdiv & T & -  & D_I  &   &      &   &      &   &      &   &        \\
	&          & -T & + & 14.50 &  &   &   &      &   &      &   &        \\
	\cline{3-7}
	&          &   &   & 14.50 & - &  D_I &   &      &   &      &   &        \\
\end{array}
\] 

So

\begin{align}
		\frac{T - D_I}{T - 14.50} = 1 + \frac{14.50 - D_I}{T-14.50} \nonumber
\end{align}

\begin{align}
	D_I = T - 0.55 \cdot (1 - 0.01 \cdot RH) (T - 14.50)
\end{align}

Assuming that:

\begin{align}
	T = \qty{38}{\degreeCelsius} \\
	RH = \qty{78}{\percent}
\end{align}

\begin{align}
	D_I &= T - 0.55 \cdot (1 - 0.01 \cdot RH) (T - 14.50) \\ 
	&= \qty{38}{\degreeCelsius} - (0.55)(1 - 0.01 \times \qty{78}{\percent} )(\qty{38}{\degreeCelsius} - 14.50) \\
	&= 35.1565
\end{align}

\begin{align}
	T &= 14.50 + \frac{14.50 - D_I}{0.55(1 - 0.01RH) - 1} \\
	&= 14.50 + \frac{14.50 - 35.1565}{0.55(1 - 0.01 \times \qty{78}{\percent} ) - 1} \\•
	&= 14.50 + \frac{-20.6565}{0.121-1} \\
	&= 14.50 + 23.5 \\
	&=\qty{38}{\degreeCelsius}
\end{align}

However, in your equation the temperature would be:

\begin{align}
	T &= \frac{D_I - (0.55)(14.5)(1 - 0.01RH)}{1 - 0.01RH} \\
	 &= \frac{35.1565 - (0.55)(14.50)(1 - 0.01 \times \qty{78}{\percent})}{1 - 0.01 \times \qty{78}{\percent}} \\
	 &= \frac{35.1565 - (7.975)(1 - 0.78)}{1 - 0.78} \\
	 &= \frac{35.1565 - 1.7545}{0.22} \\
	 &= \qty{151.83}{\degreeCelsius}
\end{align}

\begin{align}
		T &= 14.50 + \frac{14.50 - D_I}{0.55(1 - 0.01RH) - 1} \\
		&= 14.50 + \frac{14.50 - 23}{0.55(1 - 0.01 \times 75) - 1} \\
		&= 14.50 + \frac{8.5}{0.8625} \\
		&= \qty{24.35}{\degreeCelsius}
\end{align}

\begin{align}
	V_{ADC\_swing, current} = \pm \frac{\sqrt{2}I_{RMS_{max}}(R_3 + R_8)}{CT_{TURNS\_RATIO}}
\end{align}

\begin{align}
	V_{ADC\_swing, voltage} = \pm \sqrt{2}V_{RMS} \left( \frac{R_{11}}{R_{11} + R_4+R_5 + R_6 + R_7} \right)
\end{align}
\end{document}